\scanpage{52}
% L01-02
$\underline{X}=(\underline{S},\underline{A},\underline{F},\underline{R},\underline{R'},\underline{R''})$ polyèdre combinatoire $\underline{\unsure{\mathscr{R}|\mathcal{R}|R}}$ ens de ses repères.

% L03
$E$ espace affine (proj.) sur $k$ (corps)

% L04
Réalisation géom de $X$
% L05-07
\[
\left\{\begin{aligned}
\varphi_S &: \underline{S}\mathrel{\unsure{-|\longrightarrow}}E\\
\varphi_A &: \underline{A}\mathrel{\unsure{-|\longrightarrow}}\mathrm{Dr}(E)\\
\varphi_F &: \underline{F}\mathrel{\unsure{-|\longrightarrow}}\mathrm{Pl}(E)
\end{aligned}\right.
\]

% L08-09
1°) \unsure{Compatible|Compatibles} avec $R',R''$, donc $R$. On suppose de plus que

% L10-12
2°) a) Si $a\in A$, $s,s'\in S$ les deux sommets incidents $f,f'$ les deux faces incidentes, on a $\varphi_S(s)\ne\varphi_S(s')$, $\varphi_F(f)\ne\varphi_F(f')$

% L13-14
b) Si $s<f$ avec $s\in S$, $f\in F$, et si $a,a'$ sont les deux arêtes entre $s$ et $f$, on a $\varphi_A(a)\ne\varphi_A(a')$.

% L15-16
$\Rightarrow\varphi_S$ détermine $\varphi_A,\varphi_F$ / \unsure{de m.} $\varphi_F$ détermine $\varphi_A,\varphi_S$ / $\varphi_A$ détermine $\varphi_S$ donc $\varphi_F$

% L17-18, cancelled draft sentence
\deleted{Si $X$ est régulier, une telle \unsure{réalisat.|réalisation} \unsure{géom} \uline{\unsure{régulière}}}

% L19-23
\uline{Prop} Soit $\underline{r}_0=(\underline{s}_0,\underline{a}_0,\underline{f}_0)$ un repère de $X$. \unsure{Soit|Les} $\underline{s}_1$ le deuxième sommet de $\underline{a}_0$, $\underline{s}_2$ le deuxième sommet de la 2\textsuperscript{e} arête \inserted{$\unsure{a_1^+|a_0^+}$} \unsure{entre} $\underline{s}_0$ et $\underline{f}_0$, enfin $\underline{s}_3$ le deuxième sommet de la 2\textsuperscript{ième} arête \unsure{au} entre $\underline{s}_0$ et $\underline{f}_1$, la 2\textsuperscript{ième} facette \unsure{passant} par $\unsure{\underline{a}_0|\underline{f}_0}$.

% L24-25
Considérons les images $s_0,s_1,s_2,s_3$ dans $E$. Alors ces pts sont affinement indépendants
% archive sheet number: 51
